For a real nondegenerate quadratic form $Q$ of mixed signature on $\R^{n}$, $n \geq 2$, the null cone $\{Q = 0\}$ determines $Q$ up to a nonzero real scalar: in an orthonormal basis $e_{1}, \dots, e_{p}, e_{p+1}, \dots, e_{n}$, a symmetric bilinear form $S$ vanishing on the null vectors $e_{i} \pm e_{j}$ ($i \leq p < j$) has $S(e_{i}, e_{j}) = 0$ and $S(e_{i}, e_{i}) = -S(e_{j}, e_{j}) =: a$ for all such pairs, and vanishing on $3e_{i} + 4e_{k} \pm 5e_{j}$ and $5e_{i} + 3e_{j} \pm 4e_{l}$ for distinct positive $i, k$ and distinct negative $j, l$ then kills the remaining off-diagonal entries, so $S = aQ$; at $n = 2$ the first family alone suffices. For $n \geq 3$ every bijection of Minkowski space preserving the causal order is generated by Lorentz transformations, translations and dilations, so causal structure and conformal structure coincide; in $n = 2$ the two null coordinates can be reparametrised independently and the causal automorphism group is infinite-dimensional. \emph{Provenance:} the first statement is the elementary computation displayed, which the script performs on every signature up to $(3,3)$, and is the statement that the null cone fixes the conformal class \cite{HawkingEllis1973}; the null-locus criterion, that $\Cl(p,q)$ contains a null vector if and only if $pq \neq 0$, is the root paper's \cite[Prop.~3.3]{ButlerT1} (\claimref{transport.existence-locus-type-n}); the second statement is \cite{Zeeman1964}. All three enter Lemma \claimref{one-relation.third-constant-conversion-constant-between-charts} as imports.

This restates \claimref{one-relation.imported-null-cone-determines-form-zeemans-external}.
